\chapter{Catecholamines}

\section*{What Is This Test?}
Catecholamines are a group of biogenic amines that function as neurotransmitters and hormones produced by the sympathetic nervous system and the adrenal medulla. The three endogenous catecholamines are dopamine, norepinephrine (noradrenaline), and epinephrine (adrenaline). Norepinephrine is the primary neurotransmitter of post-ganglionic sympathetic neurons and is also produced by the adrenal medulla (~20\% of adrenal output). Epinephrine is produced almost exclusively by the adrenal medulla (~80\% of adrenal output), where the enzyme phenylethanolamine N-methyltransferase (PNMT) converts norepinephrine to epinephrine under the control of cortisol from the adrenal cortex. Dopamine serves as a neurotransmitter in the central nervous system and renal vasculature, and as a precursor for norepinephrine synthesis. Catecholamines are metabolized to metanephrines: norepinephrine is converted to normetanephrine by catechol-O-methyltransferase (COMT), and epinephrine is converted to metanephrine. These metanephrines are further metabolized to vanillylmandelic acid (VMA), which is excreted in the urine. Catecholamine testing (plasma free metanephrines or 24-hour urinary fractionated metanephrines and catecholamines) is primarily performed to diagnose pheochromocytoma and paraganglioma (PPGL) --- rare neuroendocrine tumors of chromaffin cells that secrete catecholamines, causing paroxysmal or sustained hypertension, palpitations, headache, and diaphoresis. Plasma free metanephrines have the highest sensitivity (96--100\%) for detecting PPGL because metanephrines are produced continuously within tumor cells (from catecholamines leaking from storage vesicles) independently of catecholamine release \cite{lenders2002}.

\section*{Alternative Names}
Plasma Free Metanephrines (Metanephrine and Normetanephrine); 24-Hour Urinary Fractionated Metanephrines; 24-Hour Urinary Catecholamines (Epinephrine, Norepinephrine, Dopamine); Plasma Catecholamines; 24-Hour Urinary VMA (Vanillylmandelic Acid); Pheochromocytoma Screen; Paraganglioma Screen; Metanephrines, Plasma Free; Catecholamine Fractionation

\section*{Principle}
Plasma free metanephrines (metanephrine and normetanephrine) are measured using liquid chromatography with electrochemical detection (LC-ECD) or liquid chromatography-tandem mass spectrometry (LC-MS/MS). LC-MS/MS is the gold standard and preferred method, offering the highest sensitivity and specificity. The plasma sample undergoes solid-phase extraction to isolate metanephrines, followed by separation on a reverse-phase C18 HPLC column and detection by tandem mass spectrometry in multiple reaction monitoring (MRM) mode. Quantification is by stable isotope dilution using deuterated internal standards (d3-metanephrine, d3-normetanephrine). Immunoassays for metanephrines exist but have lower specificity and are not recommended by Endocrine Society guidelines. Twenty-four-hour urinary fractionated metanephrines and catecholamines are also measured by LC-MS/MS or HPLC with electrochemical detection. Urinary VMA is measured by HPLC or GC-MS. The choice between plasma free metanephrines and 24-hour urinary fractionated metanephrines depends on clinical context: plasma free metanephrines have the highest sensitivity (96--100\%) but lower specificity (85--90\%), making them the preferred screening test. Urinary fractionated metanephrines have slightly lower sensitivity (90--95\%) but higher specificity (95--98\%). The combination of both tests improves diagnostic accuracy. Plasma catecholamines (epinephrine, norepinephrine) alone are NOT recommended for screening because they are released episodically and are subject to profound stress-related elevations, resulting in high false-positive rates.

\section*{History}
Epinephrine was the first hormone to be isolated and crystallized, achieved by Jokichi Takamine in 1901. Norepinephrine was identified as a sympathetic neurotransmitter by Ulf von Euler in 1946 (Nobel Prize 1970). The first case of pheochromocytoma was described pathologically by Felix Fränkel in 1886, and the first successful surgical removal was performed by César Roux in 1926 and Charles Mayo in 1927. The classic triad of headaches, palpitations, and diaphoresis as presenting symptoms was described by Melicow in 1977. Measurement of urinary catecholamines for pheochromocytoma diagnosis was established in the 1950s--1960s. The superiority of plasma metanephrines over catecholamines for pheochromocytoma detection was demonstrated by Jacques Lenders and Graeme Eisenhofer at the NIH in the 1990s, representing a major advance in diagnosis. The 2014 Endocrine Society clinical practice guideline on pheochromocytoma and paraganglioma emphasized plasma free metanephrines or 24-hour urinary fractionated metanephrines as the screening tests of choice \cite{lenders2014}. The discovery of SDHx, VHL, RET, NF1, and other susceptibility genes for PPGL has revolutionized understanding of these tumors, with 30--40\% now known to be hereditary.

\section*{Why This Test Is Ordered}
\begin{itemize}
  \item Screening for pheochromocytoma or paraganglioma in symptomatic patients with paroxysmal hypertension, palpitations, headaches, diaphoresis, and pallor --- the classic triad (headache, palpitations, sweating) occurs in <50\% of patients
  \item Evaluation of resistant hypertension (BP uncontrolled on 3 medications including a diuretic) or labile hypertension with paroxysmal exacerbations
  \item Investigation of an incidentally discovered adrenal mass (adrenal incidentaloma >1 cm) --- approximately 5--8\% of incidentalomas are pheochromocytomas
  \item Screening for PPGL in patients with a known hereditary syndrome: MEN2 (RET mutation), von Hippel-Lindau disease (VHL), neurofibromatosis type 1 (NF1), or familial paraganglioma syndromes (SDHx mutations: SDHA, SDHB, SDHC, SDHD, SDHAF2)
  \item Surveillance for recurrent or metastatic PPGL after surgical resection --- annual plasma metanephrines or 24-hour urinary metanephrines for at least 10 years
  \item Investigation of a paroxysmal event (panic attack, hypertensive crisis, unexplained cardiomyopathy, or stroke) to distinguish pheochromocytoma from anxiety or other non-catecholamine causes
  \item Evaluation of unexplained dilated cardiomyopathy or Takotsubo (stress) cardiomyopathy --- excess catecholamines can cause catecholamine-induced cardiomyopathy
  \item Detection of dopamine-secreting paraganglioma or metastatic PPGL --- plasma methoxytyramine (dopamine metabolite) elevations are characteristic
  \item Pre-operative testing before any procedure in patients with known PPGL to confirm adequate alpha-blockade (phenoxybenzamine or doxazosin)
\end{itemize}

\section*{Primary vs Secondary Test}
Plasma free metanephrines or 24-hour urinary fractionated metanephrines are primary screening tests for pheochromocytoma/paraganglioma. The Endocrine Society (2014) considers either test acceptable for initial screening, with plasma free metanephrines having the highest sensitivity and being particularly useful in high-risk populations (known mutation carriers, previous PPGL, large or suspicious adrenal mass). Plasma catecholamines alone are NOT acceptable as a screening test and are considered secondary due to high false-positive rates from stress and body position.

\section*{How This Test Is Performed}
For plasma free metanephrines: a blood sample is drawn from a vein into an EDTA (purple-top) or heparin (green-top) tube after the patient has been supine for at least 20--30 minutes in a quiet, calm environment. The specimen must be placed on ice immediately, centrifuged at 4 degrees C within 1 hour, and the plasma frozen at -20 degrees C or colder. The supine rest period is critical because upright posture and stress profoundly elevate plasma catecholamines and metanephrines. For 24-hour urinary metanephrines and catecholamines: the patient collects all urine for a complete 24-hour period in a jug containing 25 mL of 6N hydrochloric acid (HCl) as a preservative. The jug must be kept refrigerated during collection. The volume is measured and an aliquot is frozen. Urine creatinine is measured to confirm completeness of the collection (expected 15--20 mg/kg/24 hours in females, 20--25 mg/kg/24 hours in males). Urinary VMA is measured from the same 24-hour collection. Plasma catecholamines (if ordered): drawn through an indwelling catheter after at least 30 minutes of supine rest; the catheter should be placed 30 minutes before sampling to avoid the catecholamine surge from venipuncture itself.

\section*{What Preparation Is Needed}
Extensive preparation is CRITICAL for valid results. The patient must fast overnight (at least 8--10 hours, water permitted). They must avoid caffeine, nicotine, and alcohol for at least 24 hours before testing. Strenuous exercise should be avoided for 24 hours. The patient must be supine and resting for at least 20--30 minutes before blood collection in a quiet, non-stimulating environment (no TV, no phone calls). Numerous medications that interfere with catecholamine and metanephrine measurement must be discontinued if clinically safe: tricyclic antidepressants (TCAs, especially amitriptyline and imipramine, elevate norepinephrine and normetanephrine --- discontinue 7--14 days); monoamine oxidase inhibitors (MAOIs --- discontinue 14 days); levodopa/carbidopa --- discontinue 12--24 hours; sympathomimetic amines (pseudoephedrine, phenylephrine, ephedrine in decongestants and diet aids --- discontinue 48 hours); amphetamines and methylphenidate --- discontinue 48 hours; labetalol, alpha-methyldopa --- discontinue 24--48 hours; beta-blockers (particularly non-selective like propranolol, and also metoprolol to some degree --- may elevate norepinephrine and metanephrines); buspirone --- discontinue 48 hours; clonidine (if testing is to be performed, clonidine withdrawal causes a catecholamine surge --- do not suddenly discontinue); acetaminophen (paracetamol) --- can interfere with HPLC electrochemical detection and should be avoided for 48 hours; phenoxybenzamine and doxazosin --- generally continued but must be noted. The safest antihypertensives to use during testing include: calcium channel blockers (amlodipine, nifedipine), angiotensin-converting enzyme inhibitors, and angiotensin receptor blockers. The patient must inform the laboratory of ALL medications, supplements, and over-the-counter products.

\section*{Contraindications and Precautions}
Factors causing false-positive elevations and their magnitude: stress (pain, anxiety, hospitalization, venipuncture), acute illness, and surgery markedly elevate catecholamines and metanephrines; upright posture for >15--30 minutes before blood draw --- norepinephrine and normetanephrine increase 2--3 fold from supine to upright, invalidating supine reference ranges; caffeine ingestion within 24 hours (caffeine stimulates sympathetic nervous system); nicotine/tobacco use (stimulates catecholamine release); alcohol withdrawal; obstructive sleep apnea (chronic sympathetic activation, intermittent hypoxia driving catecholamine excess); congestive heart failure (compensatory sympathetic activation); severe hypertension regardless of etiology; hypoglycemia; recent myocardial infarction or stroke; intracranial pathology (subarachnoid hemorrhage, increased intracranial pressure). Drugs that cause false-positive elevations (the most clinically important): tricyclic antidepressants (TCAs --- the most common cause of false-positive metanephrine results in clinical practice because they block norepinephrine reuptake, increasing norepinephrine and normetanephrine by 2--5 fold); monoamine oxidase inhibitors (block catecholamine metabolism); levodopa (precursor of dopamine, norepinephrine, and epinephrine); sympathomimetic amines in decongestants/cold medications (pseudoephedrine, phenylephrine); amphetamines and methylphenidate; buspirone; labetalol (alpha-beta blocker, increases norepinephrine and epinephrine in urine); and alpha-methyldopa. Drugs that can cause false-negative results: metyrosine (alpha-methyl-para-tyrosine, a catecholamine synthesis inhibitor), reserpine, and ganglionic blockers. Chronic kidney disease (CKD stages 4--5): metanephrines are renally cleared and accumulate, causing mild-to-moderate elevations. Deconjugated metanephrines (total metanephrines after acid hydrolysis) are less affected. Items causing urinary interference in HPLC methods: acetaminophen, certain antibiotics, and fluorescent compounds.

\section*{Risks}
Minimal risk from venipuncture. Slight pain, bruising, or hematoma. Rare: vasovagal syncope. During 24-hour urine collection, the preservative (6N HCl) is a strong acid and MUST be handled with care to avoid skin and eye burns --- patients should be explicitly instructed not to touch the preservative. For the clonidine suppression test (rarely performed): clonidine 0.3 mg orally can cause significant hypotension, bradycardia, and sedation; the patient must be monitored for 3 hours after administration.

\section*{What Happens After the Test}
Plasma free metanephrine results are available within 3--7 days (LC-MS/MS is typically a send-out test to reference laboratories; few hospital laboratories perform this in-house). Interpretation uses established upper reference limits from large normative databases. A normal result (<upper reference limit) provides >96\% negative predictive value for excluding PPGL. Elevated results must be interpreted in context of clinical pre-test probability. Mildly elevated results (metanephrine or normetanephrine 1--2 times the upper limit) are most commonly false-positives from stress, medications, or non-PPGL conditions. Elevations 3--4 times above the upper reference limit are highly suspicious (PPV >90\%) and almost always indicate a catecholamine-secreting tumor. Elevations 2--3 times the upper limit are indeterminate and require repeat testing under optimal conditions with discontinuation of interfering medications. If repeat testing remains elevated or equivocal, the next step is abdominal/pelvic imaging: adrenal CT or MRI to localize an adrenal pheochromocytoma or extra-adrenal paraganglioma. If plasma metanephrines are clearly elevated (>4 times upper limit) and imaging identifies an adrenal mass >1 cm, the diagnosis of pheochromocytoma is established and alpha-adrenergic blockade followed by surgical resection is indicated.

\section*{Understanding Results}

\subsection*{Below Normal}
\begin{itemize}
  \item \textbf{Normal metanephrines in a patient with low pre-test probability:} Effectively excludes pheochromocytoma/paraganglioma (negative predictive value >96\%). No further evaluation for PPGL is needed.
  \item \textbf{Normal metanephrines in a patient with high pre-test probability (MEN2, SDHx mutation, prior PPGL):} Does not completely exclude a small (<1 cm) or dopamine-only secreting paraganglioma. Plasma methoxytyramine should be measured to exclude dopamine-secreting tumors. Repeat testing in 6--12 months and periodic surveillance imaging should be considered.
  \item \textbf{Normal metanephrines after surgical resection of PPGL:} Indicates complete tumor removal (biochemical cure). Annual biochemical surveillance should continue for at least 10 years because recurrence rate is 10--15\% (higher in SDHB mutation carriers, up to 30--40\%).
  \item \textbf{Low/normal metanephrines with unexplained symptoms:} If symptoms persist with normal metanephrines, consider alternative diagnoses: anxiety/panic disorder, cardiac arrhythmia (supraventricular tachycardia), migraine, menopausal hot flashes, carcinoid syndrome, mastocytosis, drug effects (sympathomimetics, cocaine), thyrotoxicosis, or hypoglycemia.
\end{itemize}

\subsection*{Above Normal}
\begin{itemize}
  \item \textbf{Elevated plasma metanephrines >3--4 times upper limit of normal:} Highly suggestive of pheochromocytoma or paraganglioma (PPV >90\%). Requires urgent localization imaging with CT or MRI of the abdomen and pelvis. Adrenal mass >1 cm confirms adrenal pheochromocytoma. No adrenal mass on CT suggests extra-adrenal paraganglioma (from sympathetic chain anywhere from skull base to pelvis). Metaiodobenzylguanidine (MIBG) I-123 scintigraphy or Ga-68 DOTATATE PET/CT should be performed for localization of paraganglioma.
  \item \textbf{Epinephrine-secreting pheochromocytoma (elevated metanephrine):} Typically adrenal in location because PNMT (converts norepinephrine to epinephrine) requires high local cortisol concentrations uniquely present in the adrenal medulla. Associated with MEN2 (RET mutations) and NF1. May present with paroxysmal episodes of palpitations, anxiety, tremor, pallor, hyperglycemia, and tachycardia.
  \item \textbf{Norepinephrine-secreting pheochromocytoma/paraganglioma (elevated normetanephrine):} May be adrenal or extra-adrenal (from sympathetic paraganglia anywhere from carotid body to organ of Zuckerkandl). Associated with von Hippel-Lindau disease (VHL) and SDHx mutations, particularly SDHB and SDHD. Patients typically have sustained or paroxysmal hypertension with headaches and diaphoresis. SDHB mutations have a 30--40\% risk of metastatic disease.
  \item \textbf{Dopamine-secreting paraganglioma (elevated methoxytyramine):} Often extra-adrenal and more likely to be malignant/metastatic. Dopamine excess may cause paradoxical normotension or hypotension (dopamine causes renal vasodilation) despite high catecholamine levels. Associated with SDHB and SDHD mutations. Methoxytyramine is the best plasma marker.
  \item \textbf{False-positive metanephrine elevations:} The most common cause is tricyclic antidepressant (TCA) use, which blocks norepinephrine reuptake and causes 2--5 fold normetanephrine elevation. Other causes: acute illness (hospitalization, ICU admission), obstructive sleep apnea (chronic sympathetic activation), recent surgery or trauma, acute myocardial infarction, congestive heart failure, severe hypertension of any etiology, hypoglycemia, alcohol withdrawal, cocaine or amphetamine use, caffeine, nicotine, and improper specimen collection with inadequate supine rest.
  \item \textbf{Equivocal metanephrine elevations (1--3 times upper limit):} Most commonly false-positives. Repeat testing after 2--4 weeks under strictly controlled conditions (optimally discontinued interfering medications, adequate supine rest). If repeat testing remains equivocal, consider the clonidine suppression test: clonidine 0.3 mg orally, measure plasma normetanephrine at baseline and 3 hours. In PPGL, normetanephrine does NOT suppress (<40\% decrease or remains above upper reference limit); in essential hypertension and normal individuals, normetanephrine suppresses >40\% and into normal range.
\end{itemize}

\section*{Normal Results}
\begin{itemize}
  \item \textbf{Plasma free metanephrine (supine):} <0.5 nmol/L (<90 pg/mL; lab-specific, varies significantly)
  \item \textbf{Plasma free normetanephrine (supine):} <0.9 nmol/L (<180 pg/mL; lab-specific)
  \item \textbf{Plasma methoxytyramine (supine):} <0.18 nmol/L (<30 pg/mL; dopamine metabolite, useful for dopamine-secreting PPGL)
  \item \textbf{24-hour urinary metanephrine:} <400 mcg/24 hours (<2 micromol/24 hours)
  \item \textbf{24-hour urinary normetanephrine:} <900 mcg/24 hours (<5 micromol/24 hours)
  \item \textbf{24-hour urinary epinephrine:} <20 mcg/24 hours (<109 nmol/24 hours)
  \item \textbf{24-hour urinary norepinephrine:} <80 mcg/24 hours (<473 nmol/24 hours)
  \item \textbf{24-hour urinary dopamine:} <400 mcg/24 hours (<2600 nmol/24 hours)
  \item \textbf{24-hour urinary VMA:} <7 mg/24 hours (<35 micromol/24 hours)
  \item \textbf{Note:} Reference intervals are method- and laboratory-specific, particularly for LC-MS/MS plasma metanephrine assays. Age-adjusted reference ranges exist (metanephrines increase slightly with age, particularly after age 60). Always use the reference interval provided by the testing laboratory.
\end{itemize}

\section*{Follow-up Tests}
\begin{itemize}
  \item \textbf{Abdominal and pelvic CT or MRI (with contrast):} First-line imaging study after biochemical confirmation of PPGL. Pheochromocytoma appears as an adrenal mass >1 cm with high attenuation on non-contrast CT (>20 HU, unlike lipid-rich adenomas), avid enhancement after IV contrast, and delayed washout. MRI: T2-weighted images show markedly hyperintense mass ("light bulb sign"), though this is not pathognomonic. Entire abdomen and pelvis must be imaged because 10--20\% of PPGLs are extra-adrenal.
  \item \textbf{I-123 or I-131 metaiodobenzylguanidine (MIBG) scintigraphy:} Functional imaging using radiolabeled MIBG, a norepinephrine analog taken up by catecholamine transporter (NET). Sensitivity 85--90\% for pheochromocytoma, 60--70\% for paraganglioma. Used to confirm catecholamine uptake, detect multiple or metastatic tumors, and determine eligibility for I-131 MIBG therapy.
  \item \textbf{Ga-68 DOTATATE PET/CT:} Somatostatin receptor imaging. Superior to MIBG for detecting paraganglioma, particularly in SDHx mutation carriers (SDHB, SDHD) where MIBG may be negative because of loss of NET expression. Increasingly the preferred functional imaging modality for PPGL.
  \item \textbf{F-18 FDG PET/CT:} Used for detecting metastatic PPGL, especially in SDHB mutation carriers where tumors are highly FDG-avid.
  \item \textbf{F-18 DOPA PET/CT:} Radiolabeled dopamine precursor taken up by catecholamine-producing cells. High sensitivity for head and neck paragangliomas and non-metastatic PPGL.
  \item \textbf{Clonidine suppression test:} For equivocal plasma normetanephrine elevations. Clonidine 0.3 mg orally suppresses norepinephrine and normetanephrine in normal and essential hypertension but NOT in PPGL. Normetanephrine decreases <40\% or remains above the upper reference limit at 3 hours: positive test (PPGL likely).
  \item \textbf{Genetic testing for susceptibility genes:} All patients with PPGL should be offered genetic counseling and testing for: RET (MEN2), VHL (von Hippel-Lindau), NF1 (neurofibromatosis type 1), SDHA, SDHB, SDHC, SDHD, SDHAF2 (familial paraganglioma syndromes), TMEM127, and MAX. Germline mutations are found in 30--40\% of apparently sporadic PPGL cases.
  \item \textbf{Plasma chromogranin A:} General neuroendocrine tumor marker that is elevated in PPGL. Not as sensitive or specific as metanephrines but can be useful for monitoring once PPGL is diagnosed.
  \item \textbf{24-hour urine creatinine:} Must be measured with every 24-hour urine collection to confirm adequacy (expected 15--25 mg/kg/24 hours). Under- or over-collection invalidates urinary metanephrine and catecholamine measurements.
\end{itemize}

\section*{References}
\bibliographystyle{unsrtnat}
\bibliography{references}

\clearpage
